\documentclass[11pt]{article}
\usepackage{fontspec}
\setmainfont{Erewhon}
\usepackage[a4paper, margin=25mm]{geometry}
\usepackage{amsmath, booktabs, graphicx, listings, verbatim, cleveref, siunitx, tikz}
% \usepackage{a4wide} and {\bf bold} and \sloppy in a comment
% A macro that marks a gap is tooling; only its uses are leftovers.
\newcommand{\TODO}[1]{\textbf{[TODO: #1]}\PackageWarning{doc}{TODO left unfilled}}
\NewDocumentCommand{\FIXME}{m}
  {\textbf{FIXME: #1}}
\def\TBD#1{[TBD: #1]}
% An input inside a definition names a file only when the macro is used.
\newcommand{\diagram}[1]{\input{figures/#1}}

\begin{document}

% Commented: {\bf x}, $$ y $$, "quoted", 3-5, 32 x 32, ..., TODO, \vspace{1em}
A line break then a comment.\\% {\bf hidden} "quoted"
A per cent sign 50\% stays text, as do ``proper'' and `single' quotes.
Ranges 3--5, a $32 \times 32$ lattice, an ellipsis \dots{} and \qty{5}{\centi\metre}.
See \cref{fig:clean} and Figure~\ref{fig:clean}, cited together~\cite{a,b}.
A label key may hold anything: \label{sec:1-2...}.
Visit \url{https://example.org/a-1-2/"x"...}.

\verb|{\bf x} "q" 3-5 ...| and \lstinline{$$ TODO $$} and \verb*!\sloppy!.

\begin{verbatim}
{\bf x} $$ y $$ "quoted" 3-5 32 x 32 ... TODO \vspace{1em}
\begin{figure}[h]
\end{verbatim}

\begin{lstlisting}[caption={Code}]
x = "string" ... 3-5 # TODO
\end{lstlisting}

\begin{comment}
\begin{figure}[H] {\bf x} TODO
\end{comment}

The phase $\phi = \sin(\omega t)$ at $T_\mathrm{max}$ and $T_{\mathrm{max}}$,
with $\operatorname{sinc} x$, $\text{min}$, $\mathrm{e}^{-x}$ and $\int f \, \mathrm{d}x$.
\begin{equation}
  W = \min\bigl(1, \mathrm{e}^{-\beta \Delta E}\bigr)
  \label{eq:max-exp}
\end{equation}

\begin{tikzpicture}
  \foreach \k in {0,...,7} {\fill (\k, 0) circle (1pt);}
  \foreach \x/\y [count=\i] in {0/1, ..., 4/5} {\node at (\x, \y) {\i};}
  \begin{axis}[xtick={0,...,10}]
  \end{axis}
\end{tikzpicture}

Standards such as ISO 80000-2 and IEC 60027-1 keep their hyphens, as do
music of the '90s and code such as \texttt{open("data.csv")}.
\begin{itemize}
  \item \cite{a} as the first word of an item.
\end{itemize}
The second figure is reached through \hyperref[fig:linked]{a named link}.

\begin{figure}[tbp]
  \centering
  \includegraphics[width=\linewidth]{plot.pdf}
  \caption{A clean figure.}
  \label{fig:clean}
\end{figure}

\begin{figure}[tbp]
  \centering
  \includegraphics[width=\linewidth]{plot.pdf}
  \caption{A figure referred to by a named link.}
  \label{fig:linked}
\end{figure}

\end{document}
